% !PRINTED 31
\weberpage{31}

Accordingly we also have the formula
\begin{equation}
 {f(x)-f(\alpha)\over x-\alpha}=Q(x),
\tag{7}
\end{equation}
where $Q(x)$ is a whole function of degree $n-1$.

If in expressions (6) we put the sign $x$ in place of the indeterminate
quantity $\alpha$, there arises a sequence of whole rational functions
of $x$, which, if for simplicity we set $a_0=1$, we write as follows:
\begin{equation}
\begin{aligned}
 f_0&=1,\\
 f_1&=x+a_1,\\
 f_2&=x^2+a_1x+a_2,\\
 &\cdots,\\
 f_{n-1}&=x^{n-1}+a_1x^{n-2}+a_2x^{n-3}+\cdots+a_{n-1}.
\end{aligned}
\tag{8}
\end{equation}
These functions $f_0,f_1,\ldots,f_{n-1}$ will later render us good
service. For the present we add the following remarks.

According to (8), the powers $1,x,x^2,\ldots,x^{n-1}$ of $x$ can be
expressed linearly in terms of the functions
$f_0,f_1,\ldots,f_{n-1}$, in such a way that only whole rational
combinations of the $a$ occur in the coefficients, for example
\[
\begin{aligned}
 1&=f_0,\\
 x&=f_1-a_1f_0,\\
 x^2&=f_2-a_1f_1+(a_1^2-a_2)f_0,\\
 &\cdots,
\end{aligned}
\]
and it follows that every whole rational function of $x$ whose degree
is not greater than $n-1$ can likewise be expressed linearly in terms
of $f_0,f_1,\ldots,f_{n-1}$ in the form
\begin{equation}
 y_0f_0+y_1f_1+\cdots+y_{n-1}f_{n-1},
\tag{9}
\end{equation}
where the coefficients $y_0,y_1,\ldots,y_{n-1}$ are independent of
$x$.

Thus, if $F(x)$ is any whole rational function of $x$, then by \S~3,
regarding $f(x)$ as divisor, one can set
\begin{equation}
 F(x)=Qf(x)+y_0f_0+y_1f_1+\cdots+y_{n-1}f_{n-1}
\tag{10}
\end{equation}
where $Q$ is also a whole rational function of $x$.

For the recursive calculation of the functions $f_r(x)$, relation (8)
gives
\begin{equation}
 f_r(x)-xf_{r-1}(x)=a_r.
\tag{11}
\end{equation}
